\section{A Structured Path Toward Productive Use}
\label{sec:path}

The design principles described in Section~\ref{sec:design_principles} can be organized
into a structured sequence for designing an AI-supported engineering workflow. The
sequence is not intended as a universal implementation lifecycle. Its purpose is to make
the main engineering questions explicit before probabilistic AI is integrated into a
productive process.

Figure~\ref{fig:controlled_ai_workflow} summarizes this processing logic. It abstracts the
architecture investigated in the underlying work into a sequence from source-grounded
information through probabilistic interpretation and explicit engineering authority to a
controlled engineering state and a verified engineering artifact.

\begin{figure*}[t]
\centering

\begin{tikzpicture}[
    font=\footnotesize,
    node distance=6mm and 7mm,
    box/.style={
        draw,
        rounded corners=1.5pt,
        align=center,
        minimum height=8mm,
        text width=34mm,
        inner sep=3pt
    },
    smallbox/.style={
        draw,
        rounded corners=1.5pt,
        align=center,
        minimum height=8mm,
        text width=29mm,
        inner sep=3pt
    },
    widebox/.style={
        draw,
        rounded corners=1.5pt,
        align=center,
        minimum height=8mm,
        text width=47mm,
        inner sep=3pt
    },
    arrow/.style={
        -{Latex[length=2mm]},
        line width=0.6pt
    },
    influence/.style={
        -{Latex[length=2mm]},
        line width=0.6pt,
        dashed
    },
    span/.style={
        draw,
        dashed,
        rounded corners=1.5pt,
        inner sep=4pt
    }
]

% ------------------------------------------------------------------
% Information basis
% ------------------------------------------------------------------

\node[smallbox] (engsource)
    {Engineering\\Sources};

\node[smallbox, right=18mm of engsource] (contextsource)
    {Context\\Sources};

\node[smallbox, right=18mm of contextsource] (processingcontext)
    {Processing and\\Orchestration Context};

\node[widebox, text width=43mm, below=7mm of engsource] (evidence)
    {Source-grounded Evidence};

\node[widebox, text width=43mm, below=7mm of contextsource] (interpretation)
    {Probabilistic Interpretation};

\draw[arrow] (engsource) -- (evidence);

\draw[arrow] (evidence.east) -- (interpretation.west);

\draw[arrow]
    (contextsource.south)
    -- (interpretation.north);

\draw[influence]
    (processingcontext.south)
    -| ($(interpretation.north east)+(-4mm,0)$);

% ------------------------------------------------------------------
% Interpretation outcome
% ------------------------------------------------------------------

\node[widebox, below=6mm of interpretation] (uncertainty)
    {Interpretation State and Decision Evidence:\\
     Consensus, Variance, Shared Ambiguity};

\node[widebox, below=6mm of uncertainty] (hea)
    {\textbf{Human Engineering Authority}\\
     Authorization of Engineering Meaning};

\node[widebox, below=6mm of hea] (authorized)
    {Authorized Engineering Information};

\draw[arrow] (interpretation) -- (uncertainty);
\draw[arrow] (uncertainty) -- (hea);
\draw[arrow] (hea) -- (authorized);

% ------------------------------------------------------------------
% Model / representation authority
% ------------------------------------------------------------------

\node[widebox, below=6mm of authorized] (modelauthority)
    {\textbf{Model Authority}\\
     Structural Placement and Representation};

\node[widebox, below=6mm of modelauthority] (controlled)
    {Controlled Engineering State\\
     \emph{e.g., Internal Engineering Model (IEM)}};

\node[widebox, below=6mm of controlled] (rules)
    {Rules, Constraints, and\\
     Target-specific Formulation};

\node[widebox, below=6mm of rules] (artifact)
    {Engineering Artifact};

\draw[arrow] (authorized) -- (modelauthority);
\draw[arrow] (modelauthority) -- (controlled);
\draw[arrow] (controlled) -- (rules);
\draw[arrow] (rules) -- (artifact);

% ------------------------------------------------------------------
% Assessment and release
% ------------------------------------------------------------------

\node[smallbox, below left=7mm and 10mm of artifact] (validation)
    {Technical\\Validation};

\node[smallbox, below right=7mm and 10mm of artifact] (review)
    {Human Model\\Review};

\coordinate (assessmentmid)
    at ($(validation.south)!0.5!(review.south)$);

\node[widebox, below=7mm of assessmentmid] (release)
    {Release};

\draw[arrow] (artifact) -- (validation);
\draw[arrow] (artifact) -- (review);

\draw[arrow] (validation.south) |- (release.west);
\draw[arrow] (review.south) |- (release.east);

% ------------------------------------------------------------------
% Cross-cutting provenance and traceability
% ------------------------------------------------------------------

\node[
    span,
    fit=(engsource)(processingcontext)(evidence)(interpretation)
        (uncertainty)(hea)(authorized)(modelauthority)
        (controlled)(rules)(artifact)(validation)(review)(release),
    label={[rotate=90,anchor=south]west:
        \textbf{Provenance and Traceability}}
] (tracebox) {};

\end{tikzpicture}

\caption{Generalized processing structure for controlled AI-supported engineering.
Engineering Sources provide the basis for source-grounded evidence, while Context Sources
support interpretation and the Processing and Orchestration Context influences how the
interpretation is produced. Probabilistic outputs remain decision evidence until explicitly
authorized. The Internal Engineering Model (IEM) illustrates one implementation of the more
general controlled engineering state.}
\label{fig:controlled_ai_workflow}
\end{figure*}


Table~\ref{tab:key_questions} condenses the corresponding design questions and their
implications for an AI-supported engineering workflow.

\begin{table*}[t]
\caption{Key Questions for Designing AI-Supported Engineering Workflows}
\label{tab:key_questions}
\centering
\renewcommand{\arraystretch}{1.15}
\begin{tabular}{p{0.12\textwidth} p{0.35\textwidth} p{0.43\textwidth}}
\toprule
\textbf{Area} &
\textbf{Key Question} &
\textbf{Design Implication} \\
\midrule

Task &
What engineering outcome is required? &
Bound the engineering task, expected result, and acceptance criteria. \\

Information &
Which information may justify an engineering statement? &
Distinguish Engineering Sources from contextual and processing information. \\

Interpretation &
Where is semantic interpretation required? &
Use probabilistic processing where explicit transformation rules are insufficient. \\

Uncertainty &
What happens when information is ambiguous, incomplete, or contradictory? &
Preserve unresolved states rather than silently completing them. \\

Authority &
Who decides whether an interpretation may be used for subsequent engineering work? &
Assign explicit engineering authority at consequential transitions. \\

Processing &
Which downstream decisions are already sufficiently determined? &
Apply explicit rules and constraints instead of repeated probabilistic interpretation. \\

Traceability &
Why does a resulting engineering state exist? &
Connect sources, interpretations, decisions, intermediate states, and artifacts. \\

Assurance &
How is the resulting artifact assessed? &
Keep generation distinct from technical validation, human review, and release. \\

\bottomrule
\end{tabular}
\end{table*}


\subsection{Define the Engineering Task and Intended Result}
\label{sec:path_task}

The starting point is the engineering task rather than the AI technology. The intended
result, the boundaries of the task, and the criteria by which the result will later be
assessed need to be defined first.

This establishes the basis for all subsequent decisions. Without a sufficiently bounded
task, it remains unclear which information is relevant, where interpretation is required,
which decisions have engineering consequences, and which form of verification is
appropriate.

The task definition should therefore clarify at least the engineering purpose, the
expected result, the scope of AI support, and the conditions under which the result can be
considered usable for subsequent engineering activities.

\subsection{Establish the Information Basis}
\label{sec:path_information}

The next step is to identify the information on which the engineering task may rely.
Available information should not only be collected but classified according to its role
within the workflow.

The underlying investigation distinguished \emph{Engineering Sources}, \emph{Context
Sources}, and the \emph{Processing and Orchestration Context}. This separation made it
possible to distinguish information that could support statements about the system from
information that supported interpretation or influenced the processing procedure.

For another engineering task, the exact classification may differ. The relevant question
remains which information may justify an engineering statement and which information is
used only to interpret, structure, or process it.

\subsection{Identify Where Interpretation Is Required}
\label{sec:path_interpretation}

Once the information basis is known, the workflow can distinguish between information that
can be processed through explicit rules and information that requires semantic
interpretation.

Probabilistic AI is particularly useful at the latter boundary. Natural-language
descriptions, incomplete structures, inconsistent terminology, and different abstraction
levels can require interpretation before they can be integrated into a controlled
engineering state.

The purpose of AI processing at this stage is not to replace the engineering decision but
to provide interpretable candidates and supporting evidence for that decision.

\subsection{Make Unresolved Information Explicit}
\label{sec:path_uncertainty}

Probabilistic interpretation may produce agreement, competing interpretations, incomplete
information, or unresolved ambiguity. These conditions should remain visible in the
processing state.

In the underlying architecture, \emph{Consensus}, \emph{Variance}, and
\emph{Shared Ambiguity} represented different outcomes of multi-perspective interpretation.
Other implementations may use different mechanisms, but the underlying requirement
remains the same: unresolved information should not disappear merely because a generative
system is able to formulate a plausible continuation.

This step therefore defines how ambiguity, contradiction, missing information, and
insufficient evidence are represented and under which conditions further processing is
blocked or escalated.

\subsection{Assign Engineering Authority}
\label{sec:path_authority}

Where AI-generated interpretations can influence subsequent engineering work, the workflow
needs an explicit decision boundary.

In the investigated architecture, the \emph{Human Engineering Authority} decided whether
an interpretation state could be used as authorized engineering information. Where a
subsequent structural representation decision was required, the \emph{Model Authority}
addressed that separate decision.

The exact role structure depends on the engineering task. The central requirement is that
the transition from a probabilistic interpretation to an authorized engineering state is
explicit and attributable to a defined authority.

Authorization should refer to a clearly identifiable state. If the underlying information
or interpretation changes materially, the corresponding decision state has changed as
well.

\subsection{Create a Controlled Engineering State}
\label{sec:path_controlled_state}

After the relevant interpretation and authority decisions have been resolved, downstream
processing should no longer depend directly on unrestricted AI output.

Instead, the authorized content can be transferred into a controlled engineering state
that provides a stable basis for subsequent processing. In the underlying implementation,
this function was realized by the \emph{Internal Engineering Model} (IEM). The specific
representation is implementation-dependent; the more general requirement is to separate
authorized engineering information from the probabilistic processing that produced its
candidate interpretation.

This boundary allows downstream processing to operate on explicit identities, relations,
rules, and approved engineering content rather than repeatedly reinterpreting the original
source material.

\subsection{Apply Rules, Constraints, and Verification}
\label{sec:path_rules}

Once the decision space has been sufficiently resolved, transformation should become
rule-bound wherever appropriate. Target-specific constraints, representation rules, and
technical validation can then be applied to the controlled engineering state.

If the required evidence, authority, or rule is missing, the unresolved condition should
remain visible. A fail-closed transition prevents the workflow from silently introducing
new engineering meaning during downstream processing.

Generation and verification remain separate concerns. The resulting artifact should be
checked using mechanisms appropriate to its purpose and representation. Depending on the
engineering task, this can include structural validation, semantic constraints, consistency
checks, tool-based validation, human review, and release criteria.

\subsection{Preserve Traceability and Integrate the Workflow}
\label{sec:path_integration}

The complete processing path should remain reconstructable from the original information
basis through interpretation and human decisions to the resulting engineering artifact.
Provenance and traceability therefore span all preceding steps rather than forming a final
documentation activity.

For multi-source workflows, this also requires preserving source-local provenance when
information is integrated into a common engineering state.

The final step toward productive use is organizational integration. Responsibilities,
review activities, quality gates, and release mechanisms need to fit the existing
engineering process. Increasing the scope of AI support also requires evaluating whether
the associated review and authorization effort remains practicable.

The resulting sequence can therefore be understood as a transition from probabilistic
interpretation to controlled engineering use: source information is interpreted where
necessary, consequential decisions remain explicitly authorized, authorized content is
transferred into a controlled state, downstream processing becomes increasingly
rule-bound, and the resulting artifact remains traceable and independently verifiable.
